\documentclass[12pt, ukrainian]{article}
\usepackage[a4paper]{geometry}
\setlength\parindent{0mm}
\geometry{top=1.1cm,bottom=1.1cm,left=1.1cm,right=1.1cm}
\pagestyle{empty}
\usepackage[T2A]{fontenc}
\usepackage[utf8]{inputenc}
\usepackage[ukrainian]{babel}
\usepackage{graphicx}
\usepackage{caption}
\usepackage{verbatim, listings, clrscode3e, fancyvrb}
\DeclareCaptionFormat{nocaption}{}
\captionsetup{format=nocaption,aboveskip=0pt,belowskip=0pt}
\usepackage[Export]{adjustbox}
\adjustboxset{max size={0.9\linewidth}{0.9\paperheight}}
\def\code#1{\Verb!#1!}
\begin{document}
{{16.05.2024 Контрольна робота, варіант  \No \textbf { 9 }} група К\_\_ ПІБ {\parbox {2.5 cm} \dotfill}\par}
{
%\begin {large}
\begin{minipage}[t]{9.7cm}
У завданнях 1--3 вказати, що буде виведено за виконання. Повідомлення інтерпретатора про помилку позначати ERROR


               
{\begin{center}Завдання {1}\end{center}}
\vspace{-4mm}

\begin{verbatim}
e = ('a','b','c','d','e',0,1,2,4)
print(e[-5])

\end{verbatim}            

\vspace{15mm}
               
{\begin{center}Завдання {2}\end{center}}
\vspace{-4mm}

\begin{verbatim}
try:
    t = {10:0, 48:0, 64:4, 35:3, 10:4}
    t[76] = 9
    for x in t.values():
        print(x, sep=' ', end=' ')
except:
    print('ERR')

\end{verbatim}            


\end{minipage}
\vline \hspace{8mm}
\begin{minipage}[t]{8.0cm}                
               
{\begin{center}Завдання {3}\end{center}}
\vspace{-4mm}

\begin{verbatim}
a = 0
n = 1

class C:
    a = 2
    
    def __init__(self):
        self.a = 3
        a = 4
        C.a = 5

obj = C()
try:
    print(a, end=' ')
    print(n, end=' ')
    print(C.a, end=' ')
    print(obj.a)
except:
    print('ERR')

\end{verbatim}            

\end{minipage}


               
{\begin{center}Завдання {4}\end{center}}
\vspace{-4mm}
Нехай змінна \kw{s} позначає послідовність цілих, по якій можна ітерувати.

Написати вираз, що перевіряє, чи правда, що всі елементи послідовності є від’ємни\-ми.\par

\vspace{10mm}
               
{\begin{center}Завдання {5}\end{center}}
\vspace{-4mm}
Змінні \kw{next} та \kw{prev} вузла списку позначають наступний та попередній за ним вузол відповідно. Починаючи з голови, у список записано послідовні цілі числа від~$-15$ до~$23$. Нехай змінна \kw{p} позначає вузол, що зберігає значення~$-3$.
Яке число зберігається у вузлі \kw{p.prev.next.next.next.prev} ?\par

               
{\begin{center}Завдання {6}\end{center}}
\vspace{-4mm}
Значенням змінної \kw{a} є цілочисловий масив типу $numpy.ndarray$ розмірів $5{\times}6$. На скільки збільшиться сума елементів масиву \kw{a} після виконання
\begin{verbatim}
a.T[-3] += 3
\end{verbatim}


Якщо код некоректний, то в якості відповіді зазначити ERROR.\par

               
{\begin{center}Завдання {10}\end{center}}
\vspace{-4mm}

Написати функцію $f$, яка для файлу, заданого шляхом, розв’язує задачу:\\ Текстовий файл містить послідовність чисел $a_1$, $a_2$, $\ldots$, $a_t$, записану у форматі одне число на рядок. Знайти середнє арифметичне чисел $a_1$, $a_4$, $a_7$, $\ldots$ (починаючи з першого через два).

В усіх випадках розміри файлів такі, що їх вміст повністю завантажений у пам'ять бути не може. Якість коду також оцінюється.



\vspace{10mm}\par

\newpage
\begin{minipage}[t]{9.7cm}
               
{\begin{center}Завдання {7}\end{center}}
\vspace{-4mm}

\begin{center}
	\adjustimage{max size={0.8\linewidth}{0.8\paperheight}}
	{../15BinTree/trees_exp/70.png}
\end{center}
Указати послідовність міток вершин дерева, зображеного на рис., що утвориться в результаті обходу дерева в прямому порядку. Мітки записувати через пробіл.\par Алгоритм починає роботу з кореня (на рис. це верхня вершина).\par

\end{minipage}
\vline \hspace{8mm}
\begin{minipage}[t]{8.0cm}
               
{\begin{center}Завдання {8}\end{center}}
\vspace{-4mm}
Для графа на рисунку з вершини 1 запущено алгоритм пошуку в глибину, що будує кістякове дерево. Списки суміжності вершин переглядаються алгоритмом за зростанням міток вершин.\par
\begin{center}
	\adjustimage{max size={0.9\linewidth}{0.9\paperheight}}
	{../16Graphs/Traversals/301.png}
\end{center}
\par Які з наведених ребер входять до складу отриманого кістякового дерева?\par
1) (5, 6)\par
2) (4, 5)\par
3) (3, 6)\par
4) (3, 5)\par
5) (2, 4)\par
6) (0, 4)\par
{\vskip 15pt} \par

\end{minipage}

               
{\begin{center}Завдання {9}\end{center}}
\vspace{-4mm}
Файл \code{'1.txt'} містить журнал з рядками, в яких через двокрапку записано поряд\-ко\-вий номер запису в журналі (ціле), код події (ціле), тип події (ERROR,\\ WARNING, INFO, STATUS), ім'я користувача (складається з латиниці), час події (фор\-мат як у прикладі), опис події.

Білі символи навколо роздільників не допускаються.

Приклад рядка файлу



\begin{verbatim}
13:404:ERROR:username:2022-05-05 13:26:something happens
\end{verbatim}


Нижче наведено код, який має в кожному рядку файли переставити місцями код та тип події.

Рядок з прикладу має бути замінений на такий



\begin{verbatim}
13:ERROR:404:username:2022-05-05 13:26:something happens
\end{verbatim}




Код:



\begin{verbatim}
regex = re.compile(r'XXX')
with open('1.txt', encoding='utf8') as f:
    lst = f.readlines()
for i in range(len(lst)):
    lst[i] = re.sub(regex, YYY, lst[i])
with open('1.txt', 'w', encoding='utf8') as f:
    f.writelines(lst)
\end{verbatim}





Який рядок треба записати на місце \kw{XXX} ?
Який рядок треба записати на місце \kw{YYY} ?\par

%\end {large}
}
%end
\end{document}

